\section{Gated Learning-Progress Exploration}
\label{sec:gated_learning_progress_exploration}

The base intrinsic score combines normalized feature-space impact and normalized region-local learning progress:
\begin{equation}
    u_t
    = \alpha_{\mathrm{impact}}\widetilde{I}_t
    + \alpha_{\mathrm{LP}}\widetilde{\mathrm{LP}}_t,
    \label{eq:base_intrinsic_score}
\end{equation}
where $\alpha_{\mathrm{impact}}\ge 0$ and $\alpha_{\mathrm{LP}}\ge 0$ determine the relative contribution of the two components.
The ungated form of the intrinsic reward is simply
\begin{equation}
    r_t^{\mathrm{int}}=u_t.
\end{equation}
This form rewards transitions that both alter the learned representation and occur in regions with recent model improvement.

The gated form adds a binary region-specific variable $g_r\in\{0,1\}$ for each visited region $r$.
The intrinsic reward becomes
\begin{equation}
    r_t^{\mathrm{int}}=g_{\rho_t}u_t.
    \label{eq:gated_intrinsic_reward}
\end{equation}
A gate value of one allows the base intrinsic score to shape the policy update, while a gate value of zero suppresses intrinsic reward for transitions assigned to that region.
The purpose of the gate is not to remove difficult regions generally, but to suppress regions that appear unproductive: high prediction error persists while local learning progress remains low.

Let $\mathcal{R}_t$ be the set of regions that have received at least one transition.
The gate uses robust global references computed from the current region statistics:
\begin{align}
    \mathrm{LP}_{\mathrm{med}}
    &= \mathrm{median}_{r\in\mathcal{R}_t}\,\mathrm{LP}(r), \\
    e_{\mathrm{med}}
    &= \mathrm{median}_{r\in\mathcal{R}_t}\,\mu_r^{\mathrm{short}}.
\end{align}
The learning-progress threshold is
\begin{equation}
    \tau_{\mathrm{LP}}=\kappa\,\mathrm{LP}_{\mathrm{med}},
\end{equation}
where $\kappa>0$ is a multiplier.
The normalized error score for region $r$ is
\begin{equation}
    s_r=\frac{\mu_r^{\mathrm{short}}}{e_{\mathrm{med}}+\varepsilon}.
\end{equation}
A region is classified as unproductive on a visit when
\begin{equation}
    \mathrm{LP}(r)<\tau_{\mathrm{LP}}
    \quad\text{and}\quad
    s_r>\tau_s,
    \label{eq:unproductive_region_condition}
\end{equation}
where $\tau_s>0$ is a high-error threshold.

The gate update includes persistence and hysteresis.
Gating is applied only after at least $R_{\min}$ regions have been visited, so that the median reference values are meaningful.
If fewer than $R_{\min}$ regions have statistics, all gates remain enabled.
When gating is active and a region with $g_r=1$ satisfies the unproductive condition for $K$ consecutive visits, the gate is set to zero.
When $g_r=0$, the gate is re-enabled only after the region satisfies
\begin{equation}
    \mathrm{LP}(r)>h\tau_{\mathrm{LP}}
\end{equation}
for two consecutive visits, with hysteresis multiplier $h>1$.
The stricter reactivation criterion reduces rapid switching caused by noise in the progress estimate.

The gated formulation treats unproductive curiosity as a local property.
A single stochastic or poorly learnable region can be suppressed without reducing intrinsic motivation everywhere else.
This differs from a purely global intrinsic-reward schedule, which can only reduce the overall strength of curiosity.
The theoretical expectation is that the ungated reward provides a low-overhead exploration signal, while the gated reward is most useful when high-error low-progress regions would otherwise attract the policy repeatedly.